% Transcription — fonds Alexandre Grothendieck, shelfmark 78, batch 3 (pages 41-60).
% Established from the facsimile archives/batches/78/batch-03.pdf (a mirror of
% https://grothendieck.umontpellier.fr/78.pdf), per the skill
% transcribe-grothendieck.
%
% Pass: Opus 5.5 (claude-opus-5-5), 2026-10-02 — first pass, unchecked against the pages by a human.
%
% Eighteen of the twenty pages are transcribed. Page 49 is a cover in his
% hand (« Cours / Polygones réguliers / II »), which becomes the section
% heading; page 53 is a blank sheet of Institut de mathématiques letterhead.
% The lower half of page 48 is blank.
%
% The batch falls into two runs.
%
%   Pages 41-48 — the end of a run on regular systems of n vertices in an
%     affine plane E over a field k (« SRnS »): the extension
%     0 → Z → Aff(S) → (Z/nZ)* → 0, the subgroup Aff±(S), the question
%     whether Z is determined by S alone, the « principe de classification
%     géométrique » (a cyclic subgroup Z ⊂ Aff(E) of order n plus an orbit S,
%     conditions 1)-3)), Prop 7, and a programme of descent: a notion of
%     épinglage, SRnS as a gerbe over the prime field, and the explicit
%     model Z_0 = μ_{n,A} ⊂ Sl(2), σ_0, s_0 = (1,1) over A = Z[1/n], whose
%     automorphism scheme is {±1}·μ_n. Page 48 ends in a boxed, crossed-out
%     construction and breaks off.
%   Pages 50-60 — « Cours Polygones réguliers II »: a plan (1. automorphismes
%     et repères des polygones combinatoires; 2. espaces affines;
%     3. métrique et formes quadratiques), a recall of the previous lecture
%     (geometric polygons, realisations, regular polygons), then combinatorial
%     polygons (S, A, R) and their flags (repères), the equivalence between
%     contours and admissible 𝔊_1-sets, 𝔊_1 = Z/2 * Z/2 = D_∞, the polygons
%     Π_n = D_n as homogeneous spaces, Corollaires 2-4 (Aut Π_n ≃ D_n),
%     § 6 on orientations, and Prop (7.) on the exact sequence
%     1 → G⁺ → G → {±1} → 1, which runs on past the batch.
%
% The run on pages 41-48 is fast and much of its connecting prose is lost;
% the course notes from page 50 on are written more carefully but remain
% abbreviated. The formulas carry better than the words throughout.
\documentclass[11pt,a4paper]{article}
\input{../preamble/grothendieck}

\folder{78}
\batch{3}
\pages{41}{60}
\dating{[à partir de 1977]}
\watermark{Édition de démonstration}
\foldertitle{[Polygones réguliers et polynômes cyclotomiques] : notes manuscrites (s.d.).}

\begin{document}

\section{Systèmes réguliers de sommets dans un plan affine}

\note{Titre de l'éditeur. La page 41 commence au milieu d'une phrase : l'argument vient du lot précédent.}

\page{41}
… fibré correspondant l'ens. de ses sections, qui est un torseur sous le \struck{groupe} \add{$\mathbb{Z}/$} des sections du fibré vectoriel correspondant, lequel est un $\mathbb{Z}/n\mathbb{Z}$-module libre de rang $1$, i.e.\ un groupe cyclique d'ordre $n$.

Si $S$ est un quasi-pol.\ combinatoire, $Z$ le groupe des rotations, le groupe de ses automorphismes \add{(noté $\mathrm{Aff}(S)$)} est \struck{canoniquement} muni d'une structure d'extension
\[
  0 \longrightarrow Z \longrightarrow \mathrm{Aff}(S) \longrightarrow (\mathbb{Z}/n\mathbb{Z})^{*} \longrightarrow 0
\]
\[
  (\mathbb{Z}/n\mathbb{Z})^{*} \simeq \mathrm{Aut}(Z)
\]
et le choix d'un $s \in S$ détermine un splitting de cette extension, permettant d'identifier $\mathrm{Aff}(S)$ au produit semi-direct $(\mathbb{Z}/n\mathbb{Z})^{*}\cdot Z$. L'image inverse du sous-groupe $\pm 1$ de $(\mathbb{Z}/n\mathbb{Z})^{*}$ dans $\mathrm{Aff}(S)$ est un sous-groupe $\mathrm{Aff}^{\pm}(S)$ de \struck{$\mathrm{Aff}(S)$} \uncertain{d'indice deux}. C'est \ill{}

\marginal{\uncertain{sous-groupe} $\mathrm{Aut}(S,\omega)$ de $\mathfrak{S}_S$, où $\omega$ est \ill{} de \uncertain{l'ensemble} des \ill{} sur $S$ compatibles avec sa str.\ \uncertain{quasi-pol.}}
\note{Note marginale écrite le long du bord gauche, lue en tournant la feuille. Sous elle, en face des deux lignes sur $\mathrm{Aff}^{\pm}(S)$, une colonne d'accolades et de signes ($\pm 1$, …) qui n'a pas été lue.}

Soit $E$ un plan affine sur un corps $k$. Un \emph{système régulier de sommets} dans $E$ est un \struck{\ill{}} couple $(S,Z)$, où $S \subset E$ est une partie finie,

\page{42}
$Z$ est une structure quasi-régulière sur $S$ (identifiée par ex.\ à un ss-groupe cyclique de $\mathfrak{S}_S$ d'ordre $n = \operatorname{card}(S)$), telle que pour une (donc toute) structure polygonale $\omega$ sur $S$ compatible avec $Z$ (i.e.\ \uncertain{liée} par $Z$), $(E,S,\omega)$ soit un polygone régulier. Je voudrais \uncertain{aussi} \ill{} que tout $u \in \mathrm{Aff}^{\pm}(S,Z)$ se prolonge de façon unique en un automorphisme affine de $E$. [NB On ne dit rien sur les autres automorphismes de $\mathrm{Aff}(S,Z)$ (il y en a bien d'autres, sauf dans le cas où $\varphi(n) = 2$\add{, ind.\ d'Euler} i.e.\ $n = $ \struck{$2$,} $3$, $4$ ou $6$).]
\marginal{Vérifier}
On voudrait aussi dire que \add{la donnée d'}\struck{un} syst.\ rég.\ de $n$ sommets dans $E$ équivaut à $2$ données~: … d'une orbite de $T$ $= (\mathbb{Z}/n\mathbb{Z})^{*}/\{\pm 1\}$ opérant sur l'ens.\ de tous les $n$-polygones réguliers dans $E$.

\page{43}
\emph{Question}. Si $(S,Z)$ est un système $n$-régulier de sommets dans $E$, $Z$ est-il déterminé de façon unique par la donnée de la seule partie $S$ de $E$~? On trouve ceci~: si $(S,\omega)$ est un $n$-polygone régulier dans $E$, le stabilisateur \add{de la partie} $S$ dans $\mathrm{Aff}(E)$ (je dis bien de $S \subset E$, non du couple $(S,\omega)$) est égal au groupe $\mathrm{Aut}(E,S,\omega)$, sauf dans le cas $\alpha = 2$ \add{ou $\alpha = -2$}~; \struck{où on trouve un groupe isomorphe} dans ces cas, identifiant $G$ à un sous-groupe de $\mathfrak{S}_S$, on trouve $\mathrm{Aff}(S,Z)$.
\marginal{vérifier~!}
D\ill{} la question \uncertain{devient} problème combinatoire, savoir~: une structure quasi-polygonale comb.\ \add{$Z$} sur $S$ est-elle déterminée par la seule connaissance de $\mathrm{Aff}^{\pm}(S) \subset \mathfrak{S}_S$, ou par celle de $\mathrm{Aff}(S) \subset \mathfrak{S}_S$~? Dans ce premier cas, $Z$ est déterminé comme le plus grand sous-groupe cyclique \emph{distingué} $\neq G = \mathrm{Aff}^{\pm}(S)$. [En effet, si $\sigma \in G \setminus Z$, on voit facilement que

\page{44}
le groupe distingué \struck{qu'il} engendré par $\sigma$ (\struck{dans} \add{égal à} \ill{}) \add{contient \ill{} de $Z$} \struck{n'est pas \ill{}}, et \ill{} est isomorphe à $\{\pm 1\}\cdot\mathbb{Z}/n'\mathbb{Z}$ (produit semi-direct) où $n' = n/2$ --- donc il n'est pas cyclique, donc $\sigma$ n'est pas contenu dans un ss-groupe invariant cyclique $\neq G$, donc tous ces groupes \add{donc} sont contenus dans $Z$, qui est le plus grand de ces groupes.] Dans le cas de $\mathrm{Aff}(S)$…
\marginal{faire dém…}
\note{Le début de la page est surchargé d'insertions et de ratures enchevêtrées~; la phrase est donnée telle qu'on peut la suivre, sans garantie sur l'ordre des ajouts.}

Donc sauf erreur, la structure quasi-polygonale combinatoire de $S \subset E$ est \add{unique} \struck{déterminée}, pour tout système régulier de $n$ sommets dans $E$.

\emph{Principe de classification géométrique}, pour $n$ fixé~: Un système régulier de $n$ sommets (\struck{en} SR$n$S) détermine
\begin{enumerate}
  \item[a)] Un sous-groupe cyclique d'ordre $n$
  \[ Z \subset \mathrm{Aff}(E) \]
  \item[b)] Une orbite $S \in Z\backslash E$
\end{enumerate}
avec les conditions suivantes
\begin{enumerate}
  \item[1)] $S$ non contenu dans une droite (cela implique que $Z$ opère \emph{fidèlement} sur $S$, qui est donc un $Z$-torseur, $Z$ étant cyclique)
  \item[2)] $\exists\, \sigma \in \mathrm{Norm}(Z)$ \struck{\ill{}} tel que $\sigma$ \struck{opère sur $Z$ par $\lambda \mapsto \lambda^{-1}$, et que $\sigma$} \add{laisse \uncertain{invariant} un pt de $S$}
\end{enumerate}

\page{45}
\struck{Réciproquement, si les conditions 1) 2) sont satisfaites, \ill{} un \ill{} de $Z$, \ill{} --- il suffit que}
\note{Les trois premières lignes de la page sont encadrées et barrées de deux traits obliques.}
$\{1,\sigma\}\cdot Z$ (produit semi-direct) opère sur $S$,… \struck{qui devient un polygone de} On a évidemment \ill{}

\emph{Prop 7}. La donnée d'un \struck{$n$-polygone} SR$n$S dans $E$ équivaut à la donnée de a) b) ci-dessus (ss-groupe \add{$Z$} cyclique d'ordre $n$ de $\mathrm{Aff}(E)$, et d'une orbite $S$ \struck{de} $Z$ dans $E$) \struck{vérifiant les conditions} \add{telles que les conditions} 1°) et 2°) ci-dessus soient vérifiées (\,$S$ non sur une droite, et condition de normalisation…) --- et éventuellement la condition minimalité~:
\begin{enumerate}
  \item[3)] $Z$ \uncertain{formé} d'ordre minimal.
\end{enumerate}

\emph{Remarque}. Suivant les cas ($p = \operatorname{car} k$, $n$), le groupe des autom.\ d'un \struck{syst.} SR$n$S est isomorphe~: $\mathrm{Aff}^{\pm}(S) \simeq D_n$, ou à $\mathrm{Aff}(S)$. On peut se poser la question, dans l'un ou l'autre cas, de \struck{trouver}

\page{46}
trouver une notion convenable d'«~épinglage~» \add{de «~repère~»} pour les SR$n$S, permettant de \struck{les} ramener la question à une structure épinglée, qui pourrait se ramener à un calcul… La théorie de descente devrait alors permettre, \struck{en définissant} comme les «~modules~» et \uncertain{le discours}, de \struck{descendre} trouver un objet épinglé canonique sur le corps premier ($\mathbb{Q}$ ou $\mathbb{F}_p$), et les autres s'en déduisent par «~torsion~» par un \struck{$(G_k)$-torseur} $G_k$-torseur \add{exceptionnel}. C'est sans doute ainsi dans le cas $G = \mathrm{Aff}(S) \simeq (\mathbb{Z}/n\mathbb{Z})^{*}\cdot(\mathbb{Z}/n\mathbb{Z})$, mais douteux dans le cas $G = \mathrm{Aff}^{\pm}(S)$. Essentiellement, la question \uncertain{revient} à celle-ci~: les SR$n$S \struck{qui} forment une gerbe sur le corps premier $\mathbb{F}_p$ (cas où $p \mid n$, on a $n = p$ \struck{premier \ill{}} ou $n = 2p$), \struck{\ill{}} \add{celle-ci est-elle \ill{}} \uncertain{une section} i.e.\ existe-t-il un SR$n$S sur le corps premier $\mathbb{F}_p$~? (Je, bien sûr, me restreins aux structures non nécessairement scindées, où \ill{} est un sous-schéma en \uncertain{fibré affine} associé à $E$, le genre sous-schéma \ill{} des points). Le mieux qu'on

\page{47}
puisse espérer dans cette direction est ceci~:
\begin{enumerate}
  \item[a)] Si $n$ \add{$\geq 3$} est premier $\neq 2$ ou de la forme $2p$, $p$ premier, trouver un SR$n$S sur $A = \mathbb{Z}[\frac{1}{n}]$
  \item[b)] Si $n$ est premier impair ou $n = 4$, trouver un SR$n$S sur $A = \mathbb{Z}$
  \item[c)] Si $n = 2p$, $p$ premier impair, trouver un SR$n$S sur $A = \mathbb{Z}[\frac{1}{2}]$.
\end{enumerate}

La solution de a) est évidente, en prenant
\[
  \begin{cases}
    Z_0 = \mu_{n,A} \subset \mathrm{Sl}(2)_{\mathbb{Z}} \subset \mathrm{Aff}\,\mathbb{E}^{2}_{/A} \\
    \sigma_0 = \begin{pmatrix} 0 & 1 \\ 1 & 0 \end{pmatrix} \in \mathrm{Sl}(2,\mathbb{Z}) \subset \mathrm{Aff}(\mathbb{E}^{2}_{A}) \\
    s_0 = (1,1) \in \mathbb{E}^{2}(A)^{\sigma} \qquad S_0 = Z \cdot s_0
  \end{cases}
\]
On trouve le \struck{groupe} SR$n$S type $(S_0,E_0)$ dont le schéma des automorphismes $G = \{\pm 1\}_A \cdot \mu_{n,A}$ (produit semi-direct) --- la catégorie fibrée des SR$n$S s'identifie donc à celle des $G$-torseurs. Ceci marche encore pour \uncertain{$\underline{A}$} $n \geq 3$, pourvu que \ill{} sur $\mathbb{Z}[1/n]$ --- donc dans les cas b) et c), on perd la car.\ $p$, où les choses sont spéciales. Bien sûr, \struck{sur le corps} \add{dans ces cas-là}, sur le corps

\note{Dans la matrice $\sigma_0$ et le stabilisateur $\mathbb{E}^{2}(A)^{\sigma}$, l'indice de $\sigma$ manque dans le manuscrit~; on le laisse tel quel.}

\page{48}
premier $\mathbb{F}_p$, il y a une \struck{représentation} \add{\uncertain{section}} de notre gerbe --- et il y a un SR$n$S scindé sur $\mathbb{F}_p$~! Mais sur $\mathbb{Z}$ (cas b)) ou $\mathbb{Z}[\frac{1}{2}]$ (cas c))~? \struck{Le plus}

\struck{\ill{} vraisemblable de regarder les cas (a) au-dessus de $\mathrm{Spec}(A[\frac{1}{p}]) = T \setminus \{p\} = T'$ ( \ill{} $T = \mathrm{Spec}(A)$ avec $\mu_{n,T'} \subset \mathrm{Sl}(2)_{T'}$, et \ill{} d'un \uncertain{point} de l'orbite}
\note{Ces lignes sont encadrées et barrées de traits obliques serrés.}

\emph{Construction}. $T = \mathrm{Spec}(A)$, \struck{\ill{}}
$E = \mathbb{E}^{2}_{T}$, $s_0$
\note{La page s'arrête ici, sa moitié inférieure blanche~: la construction n'est pas poursuivie.}

\section{Cours Polygones réguliers II}

\note{Titre de sa main, sur la couverture qui forme la page 49 («~Cours / Polygones réguliers / II~»), non transcrite comme page.}

\page{50}
\begin{enumerate}
  \item[1] Automorphismes et repères des polygones combinatoires. Polygone combinatoire-type
  \[
    \mathbb{D}_n = \{\sigma,u \mid \sigma^{2} = 1,\ \sigma u \sigma^{-1} = u^{-1},\ u^{n} = 1\}
    = \{\sigma,\sigma' \mid \sigma^{2} = \sigma'^{2} = 1,\ (\sigma'\sigma)^{n} = 1\}
  \]
  \[
    D_{\infty} = \{\sigma,u \mid \sigma^{2} = 1,\ \sigma u \sigma^{-1} = u^{-1}\}
    = \{\sigma,\sigma' \mid \sigma^{2} = \sigma'^{2} = 1\} = \mathbb{Z}_2 * \mathbb{Z}_2
  \]
  \item[2.] Espaces affines --- groupe affine $\mathrm{Aff}(E)$ ($\simeq \mathrm{Aff}(n,k) \simeq \mathrm{Gl}(n)\cdot k^{n}$)
  \item[3.] Métrique sur espaces affines et formes quadratiques. Formes quadratiques non \emph{dégénérées}.
\end{enumerate}

Rappel des points développés la fois précédente~:
\begin{enumerate}
  \item[a)] Polygone géométrique comme couple $(S,A)$
  \[
    \begin{array}{l} S \subset E \quad (\text{plan affine}) \\ A \subset \mathrm{Dr}(E) \end{array}
  \]
  Axiome~: pour la relation d'incidence habituelle ($s \in a$), $(S,A)$ est un polygone combinatoire.
  \item[b)] Réalisation géométrique d'un polygone combinatoire $\Pi = (S,A,R)$
  \[
    \left.\begin{array}{l} S \xrightarrow{\ \varphi_S\ } E \\ A \xrightarrow{\ \varphi_A\ } \mathrm{Dr}(E) \end{array}\right|
    \ \text{compatibles avec relations d'incidence}
  \]
\end{enumerate}

Si $\varphi_S(s) \neq \varphi_S(s')$ quand $s,s' \in S$ sont adjacents, \emph{alors} $\varphi_A$ est connu quand on connaît $\varphi_S$, qui peut être pris arbitrairement par ailleurs. On dira que $\varphi$ est \struck{non dégénérée} \add{fidèle} si
\[
  \varphi_S(s) \text{ incident à } \varphi_A(a) \Longrightarrow s \text{ incident à } a.
\]
Cela implique $\varphi_S$, $\varphi_A$ injectifs et $\bar{S} = \varphi_S(S)$, $\bar{A} = \varphi_A(A)$ forment un \struck{vrai} polygone dans $E$. On \struck{\ill{}} \add{dira} aussi un polygone $\Pi$-\emph{épinglé}, ou $\Pi$-\emph{polygone} (le plus souvent, $\Pi$ est fixé, $E$ et la

\page{51}
réalisation dans $E$ sont variables…)

b) Polygones réguliers. On avait proposé (dans contexte métrique, \add{sur $\mathbb{R}$})~: les déplacements \add{groupe $Z$ des} qui laissent invariant le polygone doivent être transitifs sur l'un des sommets \add{\ill{}}. Dans ce cas il est clair qu'il y a un pt invariant --- centre de gravité de l'un des sommets --- \struck{le groupe des} \add{les déplacements sont des} rotations, \struck{et le groupe} le groupe $^{R}$ des rotations de centre $o$ étant le «~groupe des angles~», canoniquement isom.\ à~:
\[
  \mathbb{R}/\mathbb{Z} \simeq U = \{z \in \mathbb{C} \mid z\bar{z} = 1\}.
\]
On trouvait, si $\zeta \in \mathbb{R}$ tel que
\[
  \begin{cases}
    \zeta^{n} = 1 \qquad \zeta^{d} \neq 1 \text{ si } 1 \leq d < n,\ d \mid n \\
    S = \{\zeta^{i} s_0 \mid i \in \mathbb{Z}/n\mathbb{Z}\} \\
    a_i = \mathrm{Dr}(s_i, s_{i+1})
  \end{cases}
\]
\note{Au-dessus de la seconde condition, une petite insertion «~$s_0 = s_1$~» précédée d'un crochet, dont la place dans la formule n'est pas claire.}
inversement, il est clair que les polygones ainsi obtenus sont «~réguliers~» au sens proposé.
\marginal{en fait, $S$ et $A$ sont des \emph{torseurs} sous $Z$, i.e.\ $Z$ est simplement transitif sur $S$ et $A$}

Si on prend les autom.\ euclidiens de $E$ sans plus \add{qui laissent $P$}, ils ont une propriété \struck{plus} de transitivité plus \emph{forte}~: \struck{ils sont} \add{leur groupe $D$ est} transitif sur l'ens.\ des \emph{repères} du polygone (et évidemment simplement transitif), puisque $S$ engendre affinement $E$.

\page{52}
Nous préférons ne pas nous donner de métrique d'avance, et prendre les autom.\ affines de $E$ qui conservent $P$, et \emph{exiger} qu'ils soient transitifs sur l'ens.\ des repères. On récupère la métrique à peu près de façon unique modulo homothétie --- mais parfois elle sera dégénérée~! (Si on exigeait dès le début une métrique non dégénérée, on perdrait des cas \struck{\ill{}} \add{\uncertain{fort intéressants}} du point de vue arithmétique --- en fait, les seuls cas vraiment nouveaux~!) Pratiquement, ce nouveau pt de vue consiste~: regarder \emph{tous} \emph{triangles} comme polygones réguliers, ainsi que \emph{tous} \emph{parallélogrammes}~! (Mais non vrais réguliers pour une métrique \emph{donnée}…)

\emph{Définition en forme}
\begin{enumerate}
  \item[1)] Polygone régulier (NB pour tt polygone, l'ens.\ des sommets engendre affinement $E$, donc un autom.\ affine est connu quand on connaît sa restriction à $S$)
  \item[2)] Réalisation géométrique régulière d'un polygone combinatoire $\Pi$…
\end{enumerate}

\page{54}
\subsection{Polygone combinatoire et repères}

(1.) Polygone combinatoire $(S,A,R)$, $S$, $A$ ens.\ \add{sommets, arêtes}, $R \subset S \times A$ \add{repères} relation, axiomes 1) $R \to A$ de degré $2$ \quad 2) $R \to S$ de degré $2$, \quad 3) Connexité. [Quand on laisse tomber l'axiome 3) de connexité, on parle de «~contour combinatoire~»~; quand on laisse tomber \struck{$R \subset$} 2) ($R \to S$ de degré $2$) on parle de graphe~; quand on \add{se donne} \struck{prend} $R \to S$ et $R \to A$, sans supposer $R \hookrightarrow S \times A$, on parle de \struck{graphe} complexe \struck{combinatoire} \add{$1$-cellulaire} combinatoire (on parlerait encore de «~graphe~» en un sens légèrement généralisé).]

(2.) Si $R \rightrightarrows \begin{smallmatrix} S \\ A \end{smallmatrix}$ est un contour, on définit
\[
  \begin{array}{ll}
    \sigma_1 & (\text{involution sans pt fixe de } R \text{ telle que } R/\sigma_1 \simeq S) \\
    \sigma_0 & (\qquad\text{---}\qquad\text{---}\qquad\text{---}\qquad R/\sigma_0 \simeq A)
  \end{array}
\]
Soit $\mathfrak{G}_1$ le groupe défini par
\[
  \mathfrak{G}_1 = \{\sigma_0,\sigma_1 \mid \sigma_0^{2} = \sigma_1^{2} = 1\} \simeq \mathbb{Z}/2\mathbb{Z} * \mathbb{Z}/2\mathbb{Z}
\]
donc on a un foncteur \ill{} de la catégorie des contours dans la cat.\ des $\mathfrak{G}_1$-ens.\ $R$ tels que
a) $\sigma_0$, $\sigma_1$ opèrent sans pt fixe dans $E$,
b) $\forall x \in E$, … $\sigma_0 x \neq \sigma_1 x$ [i.e.\ on exclut les $\mathfrak{G}_1$-orbites dans $E$ \ill{} de cardinal $2$] --- ce qui exprime que $R \to S \times A$ (où $S \overset{\mathrm{df}}{=} R/\sigma_1$, $A \overset{\mathrm{df}}{=} R/\sigma_0$) est

\page{55}
injectif \struck{\ill{}}.

\emph{Prop}. Ce foncteur
\[
  (\text{Contours comb.}) \longrightarrow \mathfrak{G}_1\text{-ens.\ ``admissibles''}
\]
est une équivalence de catégories. \add{c'est une tautologie….}

(3) Il est utile \struck{\ill{}} d'écrire $\mathfrak{G}_1$ \add{$= D_\infty$} à l'aide des générateurs \struck{$\sigma_1$ et}
\[
  \begin{cases} \sigma = \sigma_1 \\ u = \sigma_1\sigma_0 \end{cases}
  \quad\Longrightarrow\quad
  \begin{array}{l} \sigma_1 = \sigma \\ \sigma_0 = u^{-1}\sigma \end{array}
\]
\note{Dans la seconde colonne, une première écriture de $\sigma_0$ (commençant par $\sigma$) est noircie~; seule reste $u^{-1}\sigma$. Dans la première, l'expression de $u$ a été réécrite sur une rature.}
donc…
\[
  \mathfrak{G}_1 = \{\sigma,u \mid \sigma^{2} = 1,\ \sigma u \sigma^{-1} = u^{-1}\}
\]
i.e.
\[
  \mathfrak{G}_1 = \underbrace{\mathbb{Z}\cdot\mathbb{Z}/2\mathbb{Z}}_{\substack{\| \\ D_\infty\ (\text{groupe diédral infini})}}
\]
produit semi-direct avec $\sigma$ (géné.\ de $\mathbb{Z}/2\mathbb{Z}$) opérant sur $\mathbb{Z}$ par symétrie \add{les involutions}.
\marginal{NB La condition $\sigma_0 x \neq \sigma_1 x$ $\forall x \in R$ signifie qu'on exclut les «~monogones~» $=$ polyg.\ à $1$ côté…}

La condition que $\sigma_0$ et $\sigma_1$ opèrent sans pt fixe dans $R$ s'écrivent par
\[
  \left\{\begin{array}{l} \sigma x \neq x \\ u^{-1}x \neq \sigma x \end{array}\right.
  \qquad \forall x \in R
\]
La condition $\sigma_0 x \neq \sigma_1 x$ $\forall x \in R$ se transcrit par
\[
  u^{-1}x \neq x \qquad \forall x \in R
\]
donc ensemble avec les deux précédentes est équivalente à~:
\[
  \forall x \in R, \quad \operatorname{Card}\{x, \sigma x, ux\} = 3
\]
Notons que l'on a
\[
  1 \longrightarrow \mathfrak{G}_1^{\circ} \longrightarrow \mathfrak{G}_1 \longrightarrow \mathbb{Z}/2\mathbb{Z} \longrightarrow 1
\]
Notons que \ill{} est \ill{} d'ind.\ d'indice $2$, donc \ill{} splitting
\note{Cette dernière ligne est serrée contre le bord inférieur de la feuille et en grande partie illisible.}

\page{56}
(4) \emph{Cor.\ de Prop}. Les polygones comb.\ (ie contours connexes) correspondent aux espaces homogènes sous $\mathfrak{G}_1 = D_\infty$, tels que le stabilisateur d'un point quelc.\ ne contienne ni $\sigma$, ni $u$, ni $\sigma u$.

\add{Quels sont} ces sous-groupes \add{de $D_\infty$}~? \struck{\ill{}} \add{sous-groupes ne contenant} \struck{qui ne soient contenus} \add{ni $\sigma$, ni $u$}, sont réduits \add{($n \in \mathbb{N}$, $n \geq 2$)} à des sous-groupes de la forme $n\mathbb{Z}$ (ss-groupe cyclique \struck{infini} engendré par $u^{n}$) \struck{\ill{}} --- c'est donc un ss-groupe \emph{invariant} tel que le quotient soit isomorphe au groupe
\[
  \begin{aligned}
    \mathbb{D}_n &= \{\sigma,u \mid \sigma^{2} = 1,\ \sigma u \sigma^{-1} = u^{-1},\ u^{n} = 1\} \\
    &= \{\sigma_0,\sigma_1 \mid \sigma_0^{2} = \sigma_1^{2} = (\sigma_0\sigma_1)^{n} = 1\} \\
    &\simeq \mathbb{Z}_n\cdot\mathbb{Z}/2\mathbb{Z} \quad (\text{produit semi-direct},\ldots.)
  \end{aligned}
\]
où $2 \leq n \leq +\infty$ \quad ($n \in \mathbb{N}$).

(Condition évidemment suffisante. Pour la néc., il suffit de prouver que a) $H \subset \mathfrak{G}_1^{\circ} = \mathbb{Z}$ \quad b) $H \neq \mathfrak{G}_1^{\circ}$ i.e.\ $H \not\ni u$ (ce qui est vrai par hyp.~: $\sigma, u, \sigma u \notin gHg^{-1}$ $\forall g \in \mathfrak{G}_1$). Si on avait $H \not\subset \mathfrak{G}_1^{\circ}$, i.e.\ $\exists n \in \mathbb{N}$ avec $u^{n}\sigma \in H$.
\note{En b), la lettre devant l'exposant $\circ$ ressemble à $\mathbb{Z}_1$~; on lit $\mathfrak{G}_1^{\circ}$ d'après a).}

Or on a \struck{$u^{2}\sigma$}
\[
  \begin{cases}
    \operatorname{int}(u^{m})(\sigma) = u^{m}\sigma u^{-m} = u^{m}(u^{m}\sigma) = u^{2m}\sigma \\
    \operatorname{int}(u^{m})(u\sigma) = u^{2m+1}\sigma
  \end{cases}
\]
ce qui implique que, ou bien $u^{n}\sigma$ conjugué à $\sigma$ ($n$ pair), ou bien $u^{n}\sigma$ conjugué à $u\sigma$ ($n$ impair), donc $\sigma$ ou $u\sigma$ appartient à un conjugué de $H$, contrairement à l'hyp.

\page{57}
Soit $\Pi_n$ ($n \in \mathbb{N}$, $2 \leq n \leq +\infty$) le polygone combinatoire défini par $\mathbb{D}_n$, considéré comme espace homogène sous $\mathfrak{G}_1$. Alors
\[
  \Pi_n = (S_n, A_n, R_n)
\]
avec
\[
  S_n = \langle u\sigma \rangle\backslash \mathbb{D}_n \simeq \mathbb{Z}_n = \{s_i\}_{i \in \mathbb{Z}_n}
  \qquad (u\sigma = \sigma_0)
\]
\note{Devant $\backslash\mathbb{D}_n$, un premier symbole est noirci~; $u\sigma$ est écrit au-dessus, et la parenthèse «~$u\sigma$ / $\sigma_0$~» en dessous.}
[car $(u\sigma)^{2} = 1$ donc $u\sigma$ définit une section de l'ext.\ de $\mathbb{Z}_2$ par $\mathbb{Z}_n$]
\[
  A_n = \langle\sigma\rangle\backslash \mathbb{D}_n \simeq \mathbb{Z}_n = \{a_i\}_{i \in \mathbb{Z}_n}
  \qquad (\sigma = \sigma_1)
\]
enfin, la relation d'incidence est donnée par
\[
  (s_i, a_i) \quad\text{et}\quad (s_{i+1}, a_i) \quad \text{incidents.}
\]

\drawing{Une ligne brisée~: le segment $a_0$ de $s_0$ à $s_1$, puis le segment $a_1$ de $s_1$ à $s_2$, puis l'amorce d'un troisième segment, vertical, au-dessus de $s_2$.}

On retrouve la description des polyg.\ comb.\ \uncertain{d'avant}, sauf le fait que la classe d'isom.\ est caractérisée par le \ill{} des côtés $=$ \ill{} d'arêtes.

(\emph{Corollaire 2}) Si $\Pi$, $\Pi'$ sont deux polygones combinatoires iso, \struck{alors} et $r \in \mathrm{Rep}(\Pi)$, alors
\[
  \begin{array}{rcl}
    \mathrm{Isom}(\Pi,\Pi') & \xrightarrow{\ \sim\ } & \mathrm{Rep}(\Pi') \\
    \varphi & \longmapsto & \varphi(r)
  \end{array}
\]
C'est une conséquence particulière d'un \uncertain{théorème} sur les

\page{58}
espaces homogènes réguliers (ou «~distingués~»).

\emph{Exemples}~:

\emph{Corollaire 3}. Un polygone combinatoire de $n$-gone ($2 \leq n \leq \infty$) muni d'un \emph{repère} est défini à isom.\ unique près --- et est \emph{canon.\ isom.} à $\Pi_n$.

\emph{Corollaire 4}.
\[
  \begin{array}{ll}
    \mathrm{Aut}(\Pi_n) \simeq \mathbb{D}_n & (\text{canon.}) \\
    \mathrm{Aut}(\Pi) \simeq \mathbb{D}_n & (\text{non canon.}) \text{ si } \Pi \text{ $n$-gone}
  \end{array}
\]

\emph{NB}. La restriction supplémentaire $\sigma_0 x \neq \sigma_1 x$ $\forall x \in E$ (en plus de $\sigma_0$, $\sigma_1$ sans pt fixe) sur les $\mathfrak{G}_1$-ens.\ correspondants à des contours, \struck{\ill{}} revient à l'exclusion des orbites de la forme \struck{$\mathbb{Z}_2 \ldots$} $\mathbb{D}_1$ --- \uncertain{géom.}\ exclusion des polygones à $1$ côté. On pourrait les \uncertain{inclure}, en n'exigeant pas $R \hookrightarrow S \times A$.

De \ill{} façon, $=$ avec la restriction sur les orbites, il semble qu'il serait encore plus \emph{large} dans la définition, si on compare à celle des polyèdres \emph{combinatoires} --- où on exige que \add{dans} \ill{} \ill{} \ill{} \ill{} \ill{}. Je \struck{crois} \uncertain{l'exigence} \uncertain{analogue} \uncertain{reviendrait} à exiger que \ill{} \ill{} $n \geq 3$ --- \ill{}~: \ill{}

\page{59}
condition, appelons-la, que l'ens.\ \ill{} d'un polygone comme une structure sur $\underline{S}$ (ens.\ des sommets) --- à savoir un \ill{} des $2$ \ill{} circulaires --- est valable. Pour la suite, on n'aura à \uncertain{regarder} que les cas $n \geq 3$.

\subsection{6. \emph{Orientation} d'un contour}

a) Soit $R^{+} \subset R$, \struck{\ill{}} $R^{-} = R \setminus R^{+}$. Conditions équivalentes
\begin{enumerate}
  \item[a)] $\sigma_0 R^{+} = \sigma_1 R^{+} = R^{-}$
  \item[b)] $R^{+} \xrightarrow{\ \sim\ } S$, \quad $R^{+} \xrightarrow{\ \sim\ } A$
\end{enumerate}
\marginal{On dit alors que $R^{+}$ est une orientation}
c'est aussi une partie \add{$R^{+}$} de $R$ qui est le graphe d'une \emph{bijection} $S \xrightarrow{\ \sim\ } A$. Orientation \struck{\ill{}} \add{\ill{} orientations} \add{\uncertain{contours} \ill{}} \add{\ill{} $\Pi$ \ill{}}

b) \quad $\mathrm{Or}\bigl(\coprod_{i \in I} C_i\bigr) \simeq \prod \mathrm{Or}(C_i)$

\quad si $(C_i)_{i \in I}$ est une famille de contours.
\note{Sous le symbole de somme, un premier indice est noirci~; «~contours~» est écrit sous $C_i$.}

c) Si $\Pi$ est un polygone combinatoire, alors $\operatorname{card} \mathrm{Or}(\Pi) = 2$, et $R$ est \add{donc} \struck{\ill{}} disjointe de la \struck{\ill{}} famille des deux orientations \add{lesquelles sont opposées l'une de l'autre} de $\Pi$. \struck{\ill{} orientations} Les deux orientations sont les deux orbites de $\mathfrak{G}_1^{\circ}$ dans $R$.

\page{60}
Pour prouver a), notons que $\sigma_i R^{+} = R^{-}$ $\Longleftrightarrow$ $\sigma_i R^{-} = R^{+}$ ($i \in \{0,1\}$), donc si $R^{+}$ est une orientation, on aura $uR^{+} = R^{+}$ donc $\mathfrak{G}_1^{+}(R^{+}) \subset R^{+}$, et $gR^{+} = R^{-}$ si $g \in \mathfrak{G}_1 \setminus \mathfrak{G}_1^{\circ} = \mathfrak{G}_1^{-}$, \struck{Il suffit de voir que} et inversement. \ill{} dans le cas d'un polygone, \struck{\ill{}} $\simeq \Pi_n$, on voit immédiatement que $\mathfrak{G}_1^{+}$ opérant sur $R \simeq \mathbb{D}_n$ a \add{exactement} deux orbites, \add{lesquelles sont} \struck{qui} échangées par $\sigma$, savoir $\mathbb{D}_n^{+}$ et $\mathbb{D}_n^{-}$, ce sont donc ici les orientations.

(7.) \emph{Prop}. Soit $\Pi$ \add{$= (S,A,R)$} polygone combinatoire, \add{\uncertain{dièdre}}, $G = \mathrm{Aut}\,\Pi$, donc $\mathrm{Or}(\Pi)$ a deux éléments et $G$ \uncertain{opère} \ill{} de structure…
\[
  G \longrightarrow \mathfrak{S}_{\mathrm{Or}(\Pi)} \simeq \pm 1,
\]
\struck{\ill{}} Soit $G^{+}$ le noyau, \add{\ill{} ainsi} une suite exacte
\[
  1 \longrightarrow G^{+} \longrightarrow G \longrightarrow \{\pm 1\} \longrightarrow 1
\]
et la classe de $G^{-} = G \setminus G^{+}$ est d'ordre $2$, donc splitte cette suite exacte,
\begin{enumerate}
  \item[a)] \struck{$R$ est torseur} $G^{+}$ est cyclique d'ordre $n$, \struck{\ill{}} \add{\ill{}} (\struck{\ill{} $S$, $A$ \ill{}} \add{i.e.\ $A^{-1}$ de $\pm 1$}) \ill{} opère par symétrie.
  \item[b)] $R$ est un torseur sous $G$, $S$, $A$ sont torseurs sous $G^{+}$.
\end{enumerate}
\note{La proposition se poursuit vraisemblablement au-delà de ce lot.}

\end{document}
